\begin{table}[htbp]
\centering
\caption{Standardized Effect Sizes}
\label{tab:sde}
\footnotesize
\begin{tabular}{lcccccc}
\hline\hline
Outcome & $\hat{\beta}$ & SE & SD($Y$) & SDE & SE(SDE) & Class. \\
\hline
\textit{Panel A: Pooled} & & & & & & \\
\quad Log B/W hire ratio & $-$0.123 & 0.059 & 1.056 & $-$0.116 & 0.056 & Mod.\ neg. \\
\quad Log Black hires & $-$0.136 & 0.071 & 1.743 & $-$0.078 & 0.041 & Mod.\ neg. \\
\quad Log White hires & $-$0.013 & 0.022 & 0.947 & $-$0.014 & 0.023 & Small neg. \\
\quad Log B/W earnings ratio & 0.021 & 0.005 & 0.099 & 0.213 & 0.054 & Large pos. \\
\noalign{\vskip 0.5em}
\textit{Panel B: Heterogeneous} & & & & & & \\
\quad High generosity ($\geq$75\%) & 0.006 & 0.018 & 1.056 & 0.006 & 0.017 & Small pos. \\
\quad Low generosity ($<$75\%) & $-$0.138 & 0.055 & 1.056 & $-$0.131 & 0.052 & Mod.\ neg. \\
\hline\hline
\end{tabular}

\vspace{0.5em}
\begin{minipage}{\textwidth}
\scriptsize
\textit{Notes:} \textbf{Country:} United States. \textbf{Research question:} Does state-level Paid Family Leave narrow or widen the Black--White gap in new hiring from non-employment? \textbf{Policy mechanism:} PFL provides partial wage replacement (55--90\%) for workers taking family or medical leave, funded by employee payroll tax; by socializing leave costs, PFL may reduce the perceived hiring cost of workers expected to take leave, or conversely signal higher expected absence. \textbf{Outcome definition:} Log ratio of Black to White new hires from non-employment (HirA), measuring the relative flow of Black vs.\ White workers entering employment. \textbf{Treatment:} Binary; state adopts PFL benefits (8 states, 2004--2022). \textbf{Data:} QWI race microdata (Census Bureau), state $\times$ year, 2001--2023, 43 states. \textbf{Method:} Callaway--Sant'Anna staggered DiD with IPW estimation; standard errors bootstrapped (1,000 iterations); never-treated states as control group. \textbf{Sample:} Balanced state-year panel with non-missing Black and White new-hire counts; 6 treated states, 37 control states. SDE $= \hat{\beta} / \text{SD}(Y)$ where SD($Y$) is the pre-treatment standard deviation. Classification refers to magnitude, not statistical significance: Large ($|$SDE$| > 0.15$), Moderate ($0.05$--$0.15$), Small ($0.005$--$0.05$), Null ($< 0.005$).
\end{minipage}
\end{table}
